\section{Geophysical implications}

In the present section, we discuss the implications and relations of our 
theoretical results to well-established empirical approaches commonly found 
in the geophysical literature about the equivalent-layer technique.

\subsection{Application to gravity-gradient data}

The common approach to jointly process gravity-gradient data 
\cite[e.g.,][]{boggs-dransfield2004,barnes_lumley2011, martinez_li_2016} or to convert gravity disturbance data into the 
gravity-gradient tensor \cite[e.g.,][]{piauilino_etal2019_seg, oliveirajr_etal2023}, for example, 
usually presume, without any proof, the existence of a single layer of 
monopoles that satisfactorily reproduces the gravitational field and
preserves the physical relationship with the true sources.
This implies that the gravitational potential and its spatial derivatives attributed to the layer are presumed to 
accurately reproduce the corresponding field components produced by the true sources.
We provide the theoretical background for these potential-field transformations.

Let us consider, for example, the common two-step approach to convert gravity-disturbance into gravity
gradient data. The first step consists in obtaining a particular equivalent layer $\tilde{p}_{\varepsilon}$ 
by considering that the observed potential field $d^{o}$ is given by $\tfrac{\partial}{\partial z} U(\mathbf{r})$ 
(equation \ref{eq:potential-U-1st-deriv-space}) and that the kernel function 
$a(\mathbf{r}, \mathbf{r}'')$ is the vertical derivative of $\Psi$ (equation \ref{eq:Psi}).
The second step consists in using $\tilde{p}_{\varepsilon}$ to compute the gradient components.
The $uv$-component of the gradient tensor is computed by setting $D_{f} \equiv \tfrac{\partial^{2}}{\partial u \, \partial v}$
(equations \ref{eq:potential-U-2nd-deriv-space} and \ref{eq:transformed-field-continuous-layer-Neumann}),
where $u, v \in \{ x, y, z \}$.

Note that the first step consists in defining an equivalent layer $p_{\varepsilon}$ 
(equation \ref{eq:p-Fourier}) by setting $\alpha_{d} = \alpha_{a} \equiv \alpha(1, \hat{\mathbf{z}})$ 
(equations \ref{eq:observed-data-Fourier}, \ref{eq:diff-operators-generic} and \ref{eq:kernel-a-Fourier}).
By Definition \ref{def:eql-Neumann}, this particular $p_{\varepsilon}$ is a Neumann's equivalent layer,
which naturally preserves the physical relationship with the true sources.
The second step of the common approach corresponds to 
compute transformed fields $f_{\varepsilon}$ (equation \ref{eq:transformed-field-continuous-layer-Fourier}) 
defined in the Fourier domain by using $\alpha_{t} = \alpha_{f}$ (equation \ref{eq:alpha-t-generic}), with 
$\alpha_{f} \equiv \alpha(1, \hat{\mathbf{u}}) \, \alpha(1, \hat{\mathbf{v}})$ (equation \ref{eq:diff-operators-generic}), 
where the unit vectors $\hat{\mathbf{u}}$ and $\hat{\mathbf{v}}$ are given by the canonical base vectors (equation \ref{eq:canonical-base-vector}).
We can easily get from \ref{eq:diff-operators-generic} that this Fourier-domain operator $\alpha_{f}$ is equivalent to the 
space-domain operator $D_{f} \equiv \tfrac{\partial^{2}}{\partial u \, \partial v}$.
Thus, the common two-step approach to convert gravity-disturbance into gravity
gradient data via equivalent-layer technique can be theoretically explained as an application of 
a particular Neumann's equivalent layer (Definition \ref{def:eql-Neumann}).


\subsection{Classical RTP}

A predominant trend in previous studies utilizing a planar equivalent layer of dipoles for magnetic-field transformations 
is to treat the layer as a purely mathematical abstraction employed for processing total-field anomaly data, 
a perspective exemplified by the remarks of \citet{pedersen1989} concerning the research of 
\citet{bhattacharyya_chan1977}.
Relatively few studies \citep[e.g.,][]{pedersen1991, li_etal2014, reis_etal2020} have addressed the theoretical 
properties of a planar layer of dipoles capable of accurately reproducing total-field anomaly data.

Let us consider, for example, the common two-step approach to compute the RTP in the space domain, 
from total-field anomaly data, according to the pioneer work of \citet{silva1986}.
In this common approach, 
the first step consists in estimating the magnetic-moment intensity distribution on a planar layer 
with constant magnetization direction equal to that of the true sources.
The second step consists in computing the $z$-component of the magnetic induction field due 
to the estimated layer by setting its magnetization along the vertical direction.
It is usually verified a posteriori that this computed field approximates the true RTP anomaly.

We call attention to the fact that 
\citet{silva1986} cites the work of \citet{roy1962} to justify that ``the magnetic 
field due to arbitrary sources can be exactly fitted by the magnetic field due to a horizontal equivalent
layer of dipoles''. 
However, \citeauthor{roy1962} did not discuss the existence of a dipole layer that fits the observed total-field 
anomaly data and retrieves the RTP anomaly 
but presented a general outline to deduce the solutions of the Dirichlet's and Neumann's problems on a plane 
located above the sources (equation \ref{eq:U-Dirichlet-Neumann}).
In other words, the two-step approach proposed by \citet{silva1986}, which is widely accepted by the scientific community 
and we know to be correct from practice 
\citep[e.g., ][]{leao-silva_1989, guspi_novara2009, oliveirajr_etal2013, li_etal2014, reis_etal2020, takahashi-etal2022}, 
is an empirical generalization of the \citeauthor{roy1962}'s (\citeyear{roy1962}) work.
Here, we provide a theoretical explanation for this well-established approach.

Consider that (i) the observed potential field $d^{o}$ 
is satisfactorily described by the directional derivative $\hat{\mathbf{u}}^{\top} \mathbf{H}_{U}(\mathbf{r}) \, \hat{\mathbf{v}}$ 
(equation \ref{eq:potential-U-oblique-2nd-deriv-space}), where $\hat{\mathbf{u}}$ defines the constant direction of the main geomagnetic
field at the study area and $\hat{\mathbf{v}}$ the constant total-magnetization direction of the sources, and (ii)
the kernel is $a(\mathbf{r}, \mathbf{r}'') = \hat{\mathbf{u}}^{\top} \mathbf{H}_{\Psi}(\mathbf{r}, \mathbf{r}'') \, \hat{\mathbf{h}}$,
where $\mathbf{H}_{\Psi}$ is the Hessian matrix \citep[e.g.,][ p. 158]{kaplan2002} of
function $\Psi$ (equation \ref{eq:Psi}) and $\hat{\mathbf{h}}$ is a constant unit vector not necessarily equal to 
$\hat{\mathbf{v}}$.
In this case, we use $\alpha_{d} \equiv \alpha(1, \hat{\mathbf{u}}) \, \alpha(1, \hat{\mathbf{v}})$ and 
$\alpha_{a} \equiv \alpha(1, \hat{\mathbf{u}}) \, \alpha(1, \hat{\mathbf{h}})$ 
(equation \ref{eq:diff-operators-generic}) in equation \ref{eq:p-Fourier} to obtain the following equivalent layer 
in the Fourier domain:
\begin{equation}
	\mathcal{F} [ p_{\varepsilon}(\mathbf{r}'') ] =
	\begin{cases}
		k \, \tfrac{\left(\mathbf{k}^{\top} \hat{\mathbf{v}} \right)}{\left(\mathbf{k}^{\top} \hat{\mathbf{h}} \right)} \,
		\mathcal{F} \left[ U(\mathbf{r}^{\dagger}) \right] 
		e^{\, k \, (z''_{0} - z^{\dagger}_{0})} & , 
		\quad 0 < k \le k_{\varepsilon} \\
		0 & , \quad k > k_{\varepsilon} 
	\end{cases} \: .
	\label{eq:p-Fourier-dipoles}
\end{equation}
This particular equivalent layer represents a planar layer of dipoles with constant magnetization direction
$\hat{\mathbf{h}}$ and was deduced by considering
that $\hat{\mathbf{u}}$ is a unit vector having the same constant direction of the main geomagnetic field at the study area.
However, it can be easily verified that this layer can also be deduced by considering that $\hat{\mathbf{u}}$ is one of the
canonical base vectors $\hat{\mathbf{x}}$, $\hat{\mathbf{y}}$ or $\hat{\mathbf{z}}$ (equation \ref{eq:canonical-base-vector}).
In this case, the observed potential field represents a Cartesian component of the magnetic induction field.
It can also be easily verified that the same $p_{\varepsilon}$ given by equation 
\ref{eq:p-Fourier-dipoles} can be equally deduced if we consider that the 
observed potential field is the magnetic scalar potential. In this case, $\alpha_{d}$ and $\alpha_{a}$ are defined
without the term $\alpha(1, \hat{\mathbf{u}})$.

\begin{proposition}
	An equivalent layer $p_{\varepsilon}$ of dipoles defined in terms of magnetic data produced 
	by sources with constant total-magnetization direction is a Neumann's equivalent layer (Definition \ref{def:eql-Neumann}) 
	if and only if its constant magnetization direction is equal to that of the true sources.
	\label{prop:eql-Neumann-dipoles}
\end{proposition}
\begin{proof}
	In the case of a Neumann’s equivalent layer, equations \ref{eq:p-Fourier} and \ref{eq:p-Fourier-Neumann} 
	yield that $\alpha_{a} = \alpha_{d}$. 
	If the observed potential field is the scalar magnetic potential, then the corresponding
	differential operators (cf. equation \ref{eq:diff-operators-generic}) imply $\alpha_{d} \equiv \alpha(1, \hat{\mathbf{v}})$
	and $\alpha_{a} \equiv \alpha(1, \hat{\mathbf{h}})$, where $\hat{\mathbf{v}}$ and $\hat{\mathbf{h}}$ 
	are fixed unit vectors.
	In particular, $\hat{\mathbf{v}}$ is the same direction as that defining the sources
	total magnetization. The equality of these coefficients yields that 
	$(\hat{\mathbf{v}} - \hat{\mathbf{h}})^{\top}\mathbf{k} = 0$ (cf. again equation \ref{eq:diff-operators-generic}). 
	If the directions were different, i.e., if $\hat{\mathbf{h}} \ne \hat{\mathbf{v}}$, then the latter would imply that
	$\mathtt{i} k_{x}$, $\mathtt{i} k_{y}$ and $k = \sqrt{k_{x}^{2} + k_{y}^{2}}$ are linearly dependent, which is a contradiction. 
	Thus, the constant magnetization direction must be equal to that defining the total magnetization of the sources.
	
	\noindent If the observed potential field is a component of the magnetic induction field along a given
	constant direction defined by the unit vector $\hat{\mathbf{u}}$, then we have that 
	$\alpha_{d} \equiv \alpha(1, \hat{\mathbf{u}}) \alpha(1, \hat{\mathbf{v}})$ and 
	$\alpha_{a} \equiv \alpha(1, \hat{\mathbf{u}}) \alpha(1, \hat{\mathbf{h}})$. 
	The same argument as before yields that it must hold $\hat{\mathbf{h}} = \hat{\mathbf{v}}$, 
	leading to the desired conclusion.
	
	\noindent Vice-versa, assuming $\hat{\mathbf{h}} = \hat{\mathbf{v}}$ leads to $\alpha_{a} = \alpha_{d}$.
	Equations \ref{eq:p-Fourier} and \ref{eq:p-Fourier-Neumann} then imply that
	the equivalent layer under consideration must be a Neumann’s equivalent layer.
\end{proof}

We can now explain the two-step approach proposed by \citet{silva1986} from a theoretical point of view. 
The first step consists in defining a Neumann's equivalent layer $p_{\varepsilon}$ (Definition \ref{def:eql-Neumann}) in the 
Fourier domain with $\alpha_{d} = \alpha_{a} \equiv \alpha(1, \hat{\mathbf{u}}) \, \alpha(1, \hat{\mathbf{v}})$ 
(equations \ref{eq:observed-data-Fourier}, \ref{eq:diff-operators-generic} and \ref{eq:kernel-a-Fourier}). 
The second step is equivalent to 
compute the transformed potential field $f_{\varepsilon}$ (equation \ref{eq:transformed-field-continuous-layer}) defined in the Fourier domain 
by using $\alpha_{t} = \alpha_{f}$ (equation \ref{eq:alpha-t-generic}), with 
$\alpha_{f} \equiv \alpha(2, \hat{\mathbf{z}})$ (equations \ref{eq:diff-operators-generic} and \ref{eq:canonical-base-vector}).
In this case, it can be verified from equations \ref{eq:diff-operators-generic} and \ref{eq:transformed-field-continuous-layer-single-step} 
that the computed transformed field $f_{\varepsilon}$ (equation \ref{eq:transformed-field-continuous-layer}) 
approximates the RTP anomaly.

Besides, our theoretical formulation offers the theoretical link between the RTP in the space domain, 
via equivalent-layer technique, and in the Fourier domain, via classical approach \citep[e.g., ][ p. 331]{blakely1996}. 
As we have already discussed, the transformed field $f_{\varepsilon}$ computed via \ref{eq:transformed-field-continuous-layer-Fourier} 
from a previously defined equivalent layer $\tilde{p}_{\varepsilon}$ can also be directly 
computed from the observed potential field $d^{o}$ according to equation \ref{eq:transformed-field-continuous-layer-single-step}, 
in terms of the ratio $\nicefrac{\alpha_{f}}{\alpha_{d}}$.
To compute the RTP, we set $\alpha_{d} \equiv \alpha(1, \hat{\mathbf{u}}) \, \alpha(1, \hat{\mathbf{v}})$ (equations \ref{eq:observed-data-Fourier} 
and \ref{eq:diff-operators-generic}) and $\alpha_{f} \equiv \alpha(2, \hat{\mathbf{z}})$ (equation \ref{eq:diff-operators-generic}).
By using these specific $\alpha_{d}$ and $\alpha_{f}$, setting $z^{\ddagger}_{0} = z^{\dagger}_{0}$ and letting $k_{\varepsilon}$ goes to infinity 
in equation \ref{eq:transformed-field-continuous-layer-single-step}, we finally deduce the expression representing the
classical approach to compute the RTP in the Fourier domain \citep[e.g., ][ p. 331]{blakely1996}.
This theoretical link between the space- and Fourier-domain approaches to compute the RTP 
gives, for the first time, the reason by which the equivalent-layer technique ultimately suffers the same instability as 
the Fourier-domain RTP operator at low latitudes, as have already been pointed out by \citet{li_etal2014}.

PAREI AQUI

\subsection{Generalized equivalent sources}

Let us define an equivalent layer $p_{\varepsilon}$ by considering the kernel function 
$a(\mathbf{r}, \mathbf{r}'') = \Psi(\mathbf{r}, \mathbf{r}'')$ (equation \ref{eq:Psi}). 
Note that, in this case, the corresponding predicted potential field $d_{\varepsilon}$ (equation \ref{eq:predicted-data-continuous-layer}) 
represents the gravitational potential due to the equivalent layer.
This $p_{\varepsilon}$ assumes the following form in the Fourier domain:
\begin{equation}
	\mathcal{F} [ p_{\varepsilon}(\mathbf{r}'') ] =
	\begin{cases}
		k \, \alpha_{d} \, \mathcal{F} \left[ U(\mathbf{r}^{\dagger}) \right] 
		e^{\, k \, (z''_{0} - z^{\dagger}_{0})} & , 
		\quad 0 < k \le k_{\varepsilon} \\
		0 & , \quad k > k_{\varepsilon} 
	\end{cases} \: ,
	\label{eq:p-Fourier-generalized-eqs}
\end{equation}
which approximates the vertical derivative of the observed potential field on the layer plane. 
Such layers leads to potential-field transformations (equations \ref{eq:transformed-field-continuous-layer} and 
\ref{eq:transformed-field-continuous-layer-Fourier}) defined by $\alpha_{t} = \nicefrac{\alpha_{f}}{\alpha_{d}}$ (equation \ref{eq:alpha-t-generic}).

This strategy, called \textit{generalized equivalent sources}, was proposed by \citet{cordell1992} and recently adapted by 
\citet{soler_uieda2021} to perform gravity-data processing. 
Those authors considered only the particular case in which observed potential field $d^{o}$ is gravity disturbance data 
$\tfrac{\partial}{\partial z} U(\mathbf{r})$ (equation \ref{eq:observed-data-space}), 
with $\alpha_{d} \equiv \alpha(1, \hat{\mathbf{z}})$ (equations \ref{eq:observed-data-Fourier}, 
\ref{eq:diff-operators-generic} and \ref{eq:canonical-base-vector}).
In this particular case, the $p_{\varepsilon}$ given by \ref{eq:p-Fourier-generalized-eqs} 
(i) approximates $\tfrac{\partial^{2}}{\partial z^{2}} U(\mathbf{r})$ (equation \ref{eq:potential-U-2nd-deriv-space}) 
in the space domain, 
(ii) has minimum phase changes (Definition \ref{def:eql-min-phase}), with 
$\alpha_{r} = \nicefrac{1}{k}$ (equation \ref{eq:p-Fourier-min-phase}), and 
(iii) results in potential-field transformations defined by $\alpha_{t} = \nicefrac{\alpha_{f}}{k}$ (equation \ref{eq:alpha-t-generic}).

\subsection{Alternative approach to compute the RTP}

\citet{guspi_novara2009} proposed an alternative two-step approach to compute the RTP via equivalent-layer technique. 
The first step consists in obtaining an equivalent layer by considering that (i) the observed potential field $d^{o}$
is satisfactorily described by the directional derivative 
$\hat{\mathbf{u}}^{\top} \mathbf{H}_{U}(\mathbf{r}) \, \hat{\mathbf{v}}$ 
(equation \ref{eq:potential-U-oblique-2nd-deriv-space}), where $\hat{\mathbf{u}}$ defines the constant direction of the main geomagnetic
field at the study area and $\hat{\mathbf{v}}$ the constant total-magnetization direction of the sources, and (ii)
the kernel function is given by 
\begin{equation}
	a(\mathbf{r}, \mathbf{r}'') = \tfrac{\left( z''_{0} - z \right)}{2 \, \pi} \, 
	\ln{\Big[} \left( z''_{0} - z \right) + \| \mathbf{r}'' - \mathbf{r} \| {\Big]} - 
	\tfrac{\| \mathbf{r}'' - \mathbf{r} \|}{2 \, \pi} \: ,
	\label{eq:kernel-Fourier-Guspi}
\end{equation}
which represents the double integral of the function $\Psi$ (equation \ref{eq:Psi}) 
along the vertical direction $z$.
The second step of the alternative approach proposed by \citet{guspi_novara2009} computes the RTP anomaly 
by setting $f_{\varepsilon}$ (equation \ref{eq:transformed-field-continuous-layer}) with 
$t(\mathbf{r}, \mathbf{r}'') = \Psi(\mathbf{r}, \mathbf{r}'')$ (equation \ref{eq:Psi}).

Note that the first step consists in obtaining an equivalent layer $p_{\varepsilon}$ defined by 
$\alpha_{d} = (\mathbf{k}^{\top} \hat{\mathbf{u}}) (\mathbf{k}^{\top} \hat{\mathbf{v}})$ (equations \ref{eq:observed-data-Fourier} 
and \ref{eq:diff-operators-generic}) and 
$\alpha_{a} = \nicefrac{(\mathbf{k}^{\top} \hat{\mathbf{u}}) (\mathbf{k}^{\top} \hat{\mathbf{h}})}{k^{2}}$, 
where the term $\nicefrac{1}{k^{2}}$ represents the double integration along the vertical direction. In this case,  
$p_{\varepsilon}$ assumes the following form in the Fourier domain:
\begin{equation}
	\mathcal{F} [ p_{\varepsilon}(\mathbf{r}'') ] =
	\begin{cases}
		k^{3} \, \tfrac{\left(\mathbf{k}^{\top} \hat{\mathbf{v}} \right)}{\left(\mathbf{k}^{\top} \hat{\mathbf{h}} \right)} \,
		\mathcal{F} \left[ U(\mathbf{r}^{\dagger}) \right] 
		e^{\, k \, (z''_{0} - z^{\dagger}_{0})} & , 
		\quad 0 < k \le k_{\varepsilon} \\
		0 & , \quad k > k_{\varepsilon} 
	\end{cases} \: .
	\label{eq:p-Fourier-guspi}
\end{equation}
The second step consists in computing $f_{\varepsilon}$ in the Fourier domain (equation \ref{eq:transformed-field-continuous-layer-Fourier}) 
by using $\alpha_{t} = 1$ (equation \ref{eq:alpha-t-generic}).
In this case, it can be verified that the transformed potential field $\mathcal{F} [ f_{\varepsilon}(\mathbf{r}^{\ddagger}) ]$ 
(equation \ref{eq:transformed-field-continuous-layer-Fourier}) actually approximates the RTP anomaly for the particular case 
in which $\hat{\mathbf{h}} = \hat{\mathbf{v}}$, 
i.e., for the case in which the constant magnetization direction of the equivalent layer coincides with that of the true sources.
This result shows the effectiveness of our approach in providing the theoretical background not only for classical 
equivalent-layer methods, but also for unconventional ones.